\providecommand{\Ad}{\operatorname{Ad}}
\providecommand{\Hol}{\operatorname{Hol}}
\providecommand{\GL}{\operatorname{GL}}
\providecommand{\SO}{\operatorname{SO}}
\providecommand{\SU}{\operatorname{SU}}
\providecommand{\Ut}{\operatorname{U}}
\providecommand{\Ot}{\operatorname{O}}
\providecommand{\Mani}{\mathcal{M}}
\providecommand{\Lie}[1]{\mathfrak{#1}}
\providecommand{\Cl}{\operatorname{Cl}}
\providecommand{\Stab}{\operatorname{Stab}}
\providecommand{\Img}{\operatorname{Im}}
\providecommand{\rank}{\operatorname{rank}}
\providecommand{\dR}{\mathrm{dR}}
\providecommand{\softmax}{\operatorname{softmax}}

\chapter{Gauges: Gauge-Equivariant Networks}
\label{ch:gauges}

The previous chapter equipped manifolds with intrinsic geometry --- a metric,
geodesics, a connection --- and used it to transport features and build
convolutions on curved domains (\Cref{ch:geodesics}). A subtle obstruction was
left implicit there: to write down a vector, a filter, or any directional
quantity on a manifold, one must first choose a basis of the tangent space at
each point, a \emph{local reference frame} or \emph{gauge}. On a curved domain
this choice is unavoidable, it cannot be made consistently and globally, and it
is entirely arbitrary. Two practitioners analysing the same surface may pick
different frames at every point and obtain numerically different feature maps
that describe the very same geometry.

Gauge theory is the mathematical language for this freedom. Borrowed from
theoretical physics, where gauge symmetries are the internal transformations
that leave a system's equations invariant, it provides the tools to build neural
networks whose predictions are independent of the arbitrary choice of local
frame. This is the fifth and most general of the five Gs: where grids exploit a
global translation symmetry and groups a global symmetry group, gauges deal with
symmetries that act \emph{locally}, varying from point to point. The payoff is a
unifying principle. Group-equivariant networks, spherical networks, and
manifold convolutions all turn out to be special cases of gauge-equivariant
convolution on an appropriate principal bundle.

This chapter develops the framework in several movements. First, the geometric
scaffolding: principal fibre bundles, which record all possible local frames at
once, and the connections that let us compare frames at different points
(\Cref{sec:gauge:bundles,sec:gauge:transport}). Second, the symmetry itself:
gauge transformations, gauge-equivariant feature fields, and curvature as the
obstruction to a path-independent comparison
(\Cref{sec:gauge:symmetry,sec:gauge:curvature}). Third, the construction of
gauge-equivariant convolution and the icosahedral convolutional network of
\citet{cohen2019gaugeequivariantconvolutional} as a concrete, efficient
realization (\Cref{sec:gauge:conv}). Finally, we survey the algebraic and
topological generalizations the gauge viewpoint opens up --- Clifford and
geometric algebras, simplicial and cellular complexes, Hodge Laplacians,
persistent homology, and the corresponding attention architectures
(\Cref{sec:gauge:clifford,sec:gauge:topology}) --- together with positional
encodings, controlled symmetry breaking, and the message-passing synthesis that
ties the whole book together
(\Cref{sec:gauge:positional,sec:gauge:synthesis}). Throughout we lean on the
connections and parallel transport introduced in \Cref{ch:geodesics}, recasting
them in the more general bundle language that gauge symmetry demands.

\section{From global to local symmetry}
\label{sec:gauge:local}

Classical convolutional networks assume that the same transformation ---
translation --- is valid everywhere on the domain. This works for flat images
but breaks down systematically once the data carry intrinsic geometric
structure.

\begin{example}[Rotating a sphere]
\label{ex:gauge:sphere-rotation}
Consider an image of the Earth painted on a sphere and a global rotation by
$90^\circ$ about the polar axis. At the equator the rotation preserves local
shapes faithfully; near the poles the same rotation distorts them severely, and
a continent close to the north pole looks completely different afterwards. The
appropriate transformation depends on the position: a single global symmetry is
the wrong object on a curved domain.
\end{example}

On a curved domain, applying one and the same transformation everywhere can
break the natural geometric structure, introduce artificial distortions absent
from the data, and violate the physical principles governing the modelled
phenomenon. The remedy is to allow the transformation to vary from point to
point.

The physical analogy is instructive. In electromagnetism the observable fields
$\mathbf{E}$ and $\mathbf{B}$ do not depend on how we describe the
electromagnetic potential $A_\mu(x)$: we may transform
$A_\mu(x) \mapsto A_\mu(x) + \partial_\mu \Lambda(x)$ with $\Lambda$ an
arbitrary function that varies from point to point, and the fields
\[
  \mathbf{E} = -\nabla \phi - \frac{\partial \mathbf{A}}{\partial t},
  \qquad
  \mathbf{B} = \nabla \times \mathbf{A}
\]
remain invariant. This is the prototypical \emph{local} (gauge) symmetry. Just
as the laws of physics must not depend on our choice of local description,
neural representations should not depend on arbitrary choices of local reference
frame.

This suggests a natural progression in generality:
\begin{enumerate}
  \item \emph{Classical networks}: a single transformation (translation) valid
    globally;
  \item \emph{Group-equivariant networks}: a fixed finite group of
    transformations valid globally (\Cref{ch:groups});
  \item \emph{Manifold convolutions}: transformations dictated by the specific
    geometry of the manifold (\Cref{ch:geodesics});
  \item \emph{Gauge-equivariant networks}: full freedom to choose local
    transformations, subject to a consistency requirement.
\end{enumerate}

One might object: why not simply augment the data with many local
transformations? Because data augmentation hopes to achieve invariance
\emph{statistically}, across many samples and many forward passes, whereas gauge
equivariance guarantees \emph{exact} invariance by construction, in a single
forward pass. Beyond efficiency, exact equivariance prevents the network from
latching onto accidental orientations that do not reflect the true structure of
the problem.

\section{Local reference frames and principal bundles}
\label{sec:gauge:bundles}

The starting point is the structure of a homogeneous space $G/H$ studied in
\Cref{ch:foundations}. Spaces such as the sphere $S^{n-1} \cong \SO(n)/\SO(n-1)$
arise as quotients $G/H$, and identifying the stabilizer of a point $p$ with the
subgroup $H$ requires choosing a representative $g$ of the coset $gH$. Any other
representative $g' = gh$ with $h \in H$ describes the same point but gives a
different identification. This arbitrariness is the \emph{gauge freedom}, and in
general there is no canonical way to fix the choice consistently over the whole
manifold.

The mathematical solution is not to choose. Instead of picking an arbitrary
representative at each point, a \emph{principal bundle} records \emph{all}
possible choices simultaneously, attaching to each point of the base manifold
the entire group of possible reference frames.

\begin{definition}[Principal bundle]
\label{def:gauge:principal-bundle}
A \emph{principal bundle} with structure group $G$ is a quadruple
$(P, \Mani, \pi, G)$ where $P$ (the \emph{total space}) and $\Mani$ (the
\emph{base}) are smooth manifolds, $\pi\colon P \to \Mani$ is a surjective
submersion (the \emph{projection}), and $G$ is a Lie group acting \emph{freely}
on $P$ from the right --- that is, $u \cdot g = u$ for some $u$ forces $g = e$,
so $\Stab(u) = \{e\}$ for every $u$ --- subject to:
\begin{enumerate}
  \item[(a)] for each $p \in \Mani$ the fibre $\pi^{-1}(p)$ is diffeomorphic to
    $G$, the \emph{typical fibre}, so $\dim P = \dim \Mani + \dim G$; the
    diffeomorphism $\pi^{-1}(p) \cong G$ is not canonical, and fixing it amounts
    to choosing a gauge at $p$;
  \item[(b)] the action preserves fibres, $\pi(u \cdot g) = \pi(u)$; combined
    with freeness, $G$ acts freely and transitively on each fibre, making it a
    \emph{torsor} of $G$ --- a copy of $G$ with no distinguished identity;
  \item[(c)] there exist \emph{local trivializations}: every $p$ has a
    neighbourhood $U$ and a diffeomorphism $\phi_U\colon \pi^{-1}(U) \to U \times
    G$ with $\phi_U(u \cdot g) = (\pi(u), \phi_2(u)\, g)$. On overlaps,
    \emph{transition functions} $g_{UV}\colon U \cap V \to G$, defined by
    $\phi_V \circ \phi_U^{-1}(p, g) = (p, g_{UV}(p)\, g)$, relate two
    trivializations and completely encode the global topology of the bundle.
\end{enumerate}
For a systematic development see~\citep{Kobayashi1996}.
\end{definition}

\begin{remark}[Fibres as submanifolds]
\label{rem:gauge:fibres-submanifold}
Because $\pi$ is a submersion, the regular level set theorem guarantees that each
fibre $\pi^{-1}(p)$ is a smooth submanifold of $P$ of dimension
$\dim P - \dim \Mani = \dim G$. This justifies condition (a): being submanifolds
of the correct dimension on which $G$ acts freely and transitively, the
diffeomorphism $\pi^{-1}(p) \cong G$ is fixed by choosing any single point
$u_0 \in \pi^{-1}(p)$, through $g \mapsto u_0 \cdot g$.
\end{remark}

The freedom encoded in a fibre is precisely the freedom of choosing an element
within a coset $gH$. The link between the algebraic and geometric pictures is
direct.

\begin{proposition}[Homogeneous spaces are principal bundles]
\label{prop:gauge:homog-bundle}
Let $G$ be a Lie group and $H \le G$ a closed subgroup. Then
$(G, G/H, \pi, H)$, with $\pi$ the quotient projection and $H$ acting by right
multiplication, is a principal bundle.
\end{proposition}

\begin{proof}
We check the conditions of \Cref{def:gauge:principal-bundle}.
\emph{(a)} For a coset $gH \in G/H$ the fibre is $\pi^{-1}(gH) = gH$, and left
translation $L_g\colon H \to gH$, $h \mapsto gh$, is a diffeomorphism; hence each
fibre is diffeomorphic to $H$.
\emph{(b)} The right $H$-action preserves fibres: for $g \in G$, $h \in H$,
$\pi(gh) = (gh)H = gH = \pi(g)$ since $hH = H$.
\emph{(c)} Local trivializations are equivalent to local smooth sections
$s\colon U \to G$ of $\pi$, which exist for any closed subgroup $H$ of a Lie
group~\citep[Ch.~I, Prop.~5.5]{Kobayashi1996}.
\end{proof}

Thus the homogeneous spaces obtained via the orbit--stabilizer theorem of
\Cref{ch:foundations} are automatically bases of principal bundles: $S^{n-1}
\cong \SO(n)/\SO(n-1)$, $\mathbb{H}^n \cong \SO^+(n,1)/\SO(n)$, and so on. In
each case the stabilizer $H$ becomes the \emph{fibre} (the local symmetries
fixing a point of the domain), while the orbit $G/H$ is the \emph{base} on which
signals live. The gauge freedoms formalized in this chapter are precisely the
freedoms of choosing an element within each coset $gH$. \Cref{fig:gauge:bundle}
illustrates the structure.

\begin{figure}[ht]
\centering
\begin{tikzpicture}[>=Stealth, font=\small]
% Total space P (top)
\fill[gdlgreen!12, rounded corners=6pt] (-2.0, 1.0) rectangle (2.4, 4.2);
\draw[gdlgreen!50!black, dashed, rounded corners=6pt] (-2.0, 1.0) rectangle (2.4, 4.2);
\node[gdlgreen!50!black, anchor=south] at (0.2, 4.2) {$\pi^{-1}(U)$};
\draw[thick, gdlblue!40!black, rounded corners=8pt, dashed] (-4.4, 0.8) rectangle (3.0, 4.7);
\node[gdlblue!40!black, anchor=south] at (-0.7, 4.7) {$P$};
% fibre outside U
\draw[thick, gdlorange!60, fill=gdlorange!15] (-3.4, 2.6) ellipse (0.28cm and 0.75cm);
\foreach \yy/\op in {3.1/65, 2.6/50, 2.1/35} \fill[gdlred!\op] (-3.4, \yy) circle (1.5pt);
\draw[gdlgray!40, dashed] (-3.4, 1.85) -- (-3.4, -0.2);
\node[gdlorange!70!black, left] at (-3.7, 2.6) {$G$};
% fibres over U
\foreach \xx in {-1.2, 0.2, 1.6} {
  \draw[thick, gdlorange!70, fill=gdlorange!18] (\xx, 2.6) ellipse (0.28cm and 0.75cm);
  \foreach \yy/\op in {3.1/65, 2.6/50, 2.1/35} \fill[gdlred!\op] (\xx, \yy) circle (1.5pt);
  \draw[gdlgray!40, dashed] (\xx, 1.85) -- (\xx, -0.2);
}
\node[gdlorange!70!black, right] at (1.95, 3.2) {$\pi^{-1}(p)\cong G$};
% projection
\draw[->, very thick, gdlpurple!70] (0, 0.7) -- node[right] {$\pi$} (0, 0.2);
% base M
\fill[gdlblue!15] (-4.4, -0.6) .. controls (-3, 0.0) and (-1, -0.8) .. (0.5, -0.5)
   .. controls (2, -0.2) and (2.4, -0.7) .. (3.0, -0.55)
   -- (3.0, -1.4) .. controls (2.4, -1.5) and (2, -1.2) .. (0.5, -1.5)
   .. controls (-1, -1.8) and (-3, -1.0) .. (-4.4, -1.6) -- cycle;
\fill[gdlgreen!22, rounded corners=4pt] (-2.0, -0.55) rectangle (2.4, -1.35);
\node[gdlgreen!50!black] at (0.2, -0.95) {$U$};
\node[gdlblue!70!black] at (-3.7, -0.3) {$\Mani$};
\foreach \xx in {-1.2, 0.2, 1.6} \fill[gdlgreen!50!black] (\xx, -0.75) circle (1.5pt);
\node[gdlgreen!50!black, below, font=\tiny] at (0.2, -0.85) {$p$};
\fill[gdlblue!60] (-3.4, -0.75) circle (1.5pt);
% product side
\begin{scope}[shift={(7.0,0.0)}]
  \fill[gdlgreen!15] (-1.0, 0.0) rectangle (2.0, 3.2);
  \draw[thick, gdlgreen!50!black] (-1.0, 0.0) rectangle (2.0, 3.2);
  \draw[->, thick] (-1.0, 0.0) -- (2.3, 0.0); \node[below] at (0.6,-0.05) {$U$};
  \draw[->, thick] (-1.0, 0.0) -- (-1.0, 3.5); \node[left] at (-1.0,1.6) {$G$};
  \node[gdlgreen!40!black] at (0.6, 3.5) {$U \times G$};
  \foreach \xpos in {-0.3, 0.6, 1.5} {
    \draw[gdlorange!60, thick] (\xpos, 0.1) -- (\xpos, 3.1);
    \foreach \yy/\op in {2.4/65,1.6/50,0.8/35} \fill[gdlred!\op] (\xpos, \yy) circle (1.5pt);
  }
\end{scope}
\draw[<->, thick, gdlpurple!80] (2.6, 2.6) -- node[above] {$\phi_U$} node[below, font=\scriptsize, gdlpurple!60] {$\cong$} (5.9, 1.6);
\end{tikzpicture}
\caption[Principal bundle]{A principal bundle and its local trivialization. The
total space $P$ projects onto the base manifold $\Mani$ through $\pi$; each fibre
$\pi^{-1}(p)$ is diffeomorphic to the structure group $G$ and collects all
possible local frames at $p$. Over a neighbourhood $U$ the diffeomorphism
$\phi_U$ presents the bundle as the product $U \times G$, but its global topology
may be twisted.}
\label{fig:gauge:bundle}
\end{figure}

\begin{example}[Orthonormal frame bundle]
\label{ex:gauge:frame-bundle}
Let $(\Mani, g)$ be an $n$-dimensional Riemannian manifold. Its
\emph{orthonormal frame bundle} $\mathrm{FM}(\Mani)$ is the principal bundle with
structure group $\Ot(n)$ whose total space consists of all orthonormal frames
$(p, e_1, \ldots, e_n)$ with $\{e_i\}$ an orthonormal basis of $T_p\Mani$; the
projection $(p, e_1, \ldots, e_n) \mapsto p$ forgets the frame, and $A \in \Ot(n)$
acts by $(p, e_1, \ldots, e_n) \cdot A = (p, \sum_j A_{j1} e_j, \ldots, \sum_j
A_{jn} e_j)$. Each fibre collects every possible choice of orthonormal frame at
$p$ --- precisely the local gauge freedom. For a surface ($n = 2$) the relevant
rotational freedom is $\SO(2)$.
\end{example}

\begin{example}[Trivial bundle]
\label{ex:gauge:trivial-bundle}
The simplest bundle is the product $P = \Mani \times G$ with $\pi(p, g) = p$ and
action $(p, h) \cdot g = (p, hg)$. It admits a global section $s(p) = (p, e)$ ---
a globally consistent choice of gauge, with $e$ the identity. Not every bundle
admits a global section; whether it does is a topological question, taken up in
\Cref{sec:gauge:curvature}.
\end{example}

In the language of geometric deep learning, the base $\Mani$ is where the data
live (for instance the sphere $S^2$ for spherical signals), while the structure
group $G$ encodes the relevant local symmetries ($\SO(2)$ for rotations of a
tangent frame). The fibre over each point contains all possible local
orientations or reference frames there.

\section{Connections and parallel transport}
\label{sec:gauge:transport}

A connection on a principal bundle prescribes how to transport information
between distinct fibres in a way compatible with the group structure. It splits
the tangent space of the total space into a \emph{vertical} part along the
fibre and a \emph{horizontal} part transverse to it. This comparison is what
makes operations such as convolution definable on non-Euclidean domains.

\begin{definition}[Vertical space and fundamental vector field]
\label{def:gauge:vertical}
For $u \in P$ the \emph{vertical space} is
$V_u P = \ker(\mathrm{d}\pi_u) = \{X \in T_u P : \mathrm{d}\pi_u(X) = 0\}$, the
directions tangent to the fibre. For $A \in \Lie{g}$, the Lie algebra of $G$, the
\emph{fundamental vector field} generated by $A$ at $u$ is
\[
  A^*_u = \left.\frac{\mathrm{d}}{\mathrm{d}t}\right|_{t=0} u \cdot \exp(tA),
\]
the velocity of the orbit curve $t \mapsto u \cdot \exp(tA)$. Since this curve
stays in the fibre $\pi^{-1}(\pi(u))$, we have $A^*_u \in V_u P$.
\end{definition}

\begin{proposition}[Canonical isomorphism $V_u P \cong \Lie{g}$]
\label{prop:gauge:vertical-iso}
For each $u \in P$ the map $\sigma_u\colon \Lie{g} \to V_u P$, $\sigma_u(A) =
A^*_u$, is a linear isomorphism.
\end{proposition}

\begin{proof}
Consider the orbit map $\Phi_u\colon G \to P$, $\Phi_u(g) = u \cdot g$. Its
differential at the identity is $\sigma_u = (\mathrm{d}\Phi_u)_e\colon T_e G =
\Lie{g} \to T_u P$, with image inside $V_u P$ by
\Cref{def:gauge:vertical}; being the differential of a smooth map it is linear.
For injectivity, freeness of the action makes $\Phi_u$ injective and an
immersion: if $(\mathrm{d}\Phi_u)_e(A) = 0$ the curve $t \mapsto u \cdot \exp(tA)$
has zero initial velocity, which by freeness forces $A = 0$. For surjectivity,
since $\pi$ is a submersion, $\dim V_u P = \dim P - \dim \Mani = \dim G = \dim
\Lie{g}$, and an injective linear map between spaces of equal finite dimension is
an isomorphism.
\end{proof}

\begin{definition}[Connection]
\label{def:gauge:connection}
A \emph{connection} on $(P, \Mani, \pi, G)$ is a Lie-algebra-valued $1$-form
$\omega \in \Omega^1(P, \Lie{g})$ such that
\begin{enumerate}
  \item \emph{normalization}: $\omega(A^*_u) = A$ for all $A \in \Lie{g}$ and
    $u \in P$;
  \item \emph{equivariance}: $R_g^* \omega = \Ad_{g^{-1}} \omega$ for all
    $g \in G$, where $R_g$ is the right action and $\Ad$ the adjoint
    representation.
\end{enumerate}
Equivalently, a connection defines a \emph{horizontal distribution}
$H_u P = \ker(\omega_u) \subset T_u P$ complementary to the vertical space, so
that $T_u P = H_u P \oplus V_u P$ and $H_{u \cdot g} P = (R_g)_* H_u P$.
\end{definition}

The horizontal distribution specifies which directions in the total space are
``purely spatial'' (parallel to the base) versus ``purely gauge'' (along the
fibre). A curve $\tilde\gamma$ in $P$ is \emph{horizontal} if
$\dot{\tilde\gamma}(t) \in H_{\tilde\gamma(t)} P$ for all $t$; such curves
represent transport of the frame ``without rotation.'' This is exactly the
bundle-theoretic form of the affine connection of \Cref{ch:geodesics}.

\begin{definition}[Parallel transport]
\label{def:gauge:parallel-transport}
Given a connection and a smooth path $\gamma\colon [0,1] \to \Mani$,
\emph{parallel transport} is the map
$\Gamma_\gamma\colon \pi^{-1}(\gamma(0)) \to \pi^{-1}(\gamma(1))$ defined as
follows: for $u_0 \in \pi^{-1}(\gamma(0))$ there is a unique \emph{horizontal
lift} $\tilde\gamma\colon [0,1] \to P$ with (i) $\pi \circ \tilde\gamma =
\gamma$, (ii) $\tilde\gamma(0) = u_0$, and (iii) $\dot{\tilde\gamma}(t) \in
H_{\tilde\gamma(t)} P$ for all $t$; then $\Gamma_\gamma(u_0) = \tilde\gamma(1)$.
\end{definition}

\begin{theorem}[Properties of parallel transport]
\label{thm:gauge:transport-props}
Parallel transport satisfies:
\begin{enumerate}
  \item \emph{equivariance}: $\Gamma_\gamma(u \cdot g) = \Gamma_\gamma(u) \cdot
    g$ for all $g \in G$;
  \item \emph{composition}: $\Gamma_{\gamma_2 * \gamma_1} = \Gamma_{\gamma_2}
    \circ \Gamma_{\gamma_1}$ for concatenated paths;
  \item \emph{inversion}: $\Gamma_{\gamma^{-1}} = (\Gamma_\gamma)^{-1}$, where
    $\gamma^{-1}(t) = \gamma(1-t)$.
\end{enumerate}
\end{theorem}

\begin{proof}
\emph{Equivariance.} Let $\tilde\gamma$ be the horizontal lift from $u$ and set
$\tilde\gamma_g(t) = \tilde\gamma(t) \cdot g$. Then $\pi(\tilde\gamma_g(t)) =
\pi(\tilde\gamma(t)) = \gamma(t)$; $\tilde\gamma_g(0) = u \cdot g$; and
$\dot{\tilde\gamma}_g(t) = (R_g)_* \dot{\tilde\gamma}(t) \in (R_g)_*
H_{\tilde\gamma(t)} P = H_{\tilde\gamma(t)\cdot g} P$ by the $G$-invariance of the
horizontal distribution. By uniqueness of the lift, $\tilde\gamma_g$ is the lift
from $u \cdot g$, so $\Gamma_\gamma(u \cdot g) = \tilde\gamma(1)\cdot g =
\Gamma_\gamma(u)\cdot g$.

\emph{Composition.} Let $\tilde\gamma_1$ be the horizontal lift of $\gamma_1$
from $u_0$, ending at $u_1 = \tilde\gamma_1(1)$, and $\tilde\gamma_2$ the lift of
$\gamma_2$ from $u_1$. The concatenation $\tilde\gamma$ projects to
$\gamma_2 * \gamma_1$, starts at $u_0$, and is horizontal on each piece; by
uniqueness it is the lift of $\gamma_2 * \gamma_1$ from $u_0$, with endpoint
$\tilde\gamma_2(1) = \Gamma_{\gamma_2}(\Gamma_{\gamma_1}(u_0))$.

\emph{Inversion.} With $\tilde\gamma^-(t) = \tilde\gamma(1-t)$ we have
$\pi(\tilde\gamma^-(t)) = \gamma^{-1}(t)$, $\tilde\gamma^-(0) = u_1$, and
$\dot{\tilde\gamma}^-(t) = -\dot{\tilde\gamma}(1-t) \in H_{\tilde\gamma(1-t)} P$,
since horizontal subspaces are vector spaces. By uniqueness this is the lift of
$\gamma^{-1}$ from $u_1$, giving $\Gamma_{\gamma^{-1}}(u_1) = u_0 =
\Gamma_\gamma^{-1}(u_1)$~\citep[Ch.~II, Prop.~3.2]{Kobayashi1996}.
\end{proof}

Working in a local section, the horizontal lift satisfies
$\omega(\dot{\tilde\gamma}) = 0$, which expressed through a group-valued curve
$g(t)$ becomes the linear ordinary differential equation
\[
  \frac{\mathrm{d}g}{\mathrm{d}t} + \omega(\dot\gamma(t))\, g(t) = 0,
\]
with the path-ordered exponential
$g(1) = \mathcal{P}\exp\!\big(-\!\int_0^1 \omega(\dot\gamma)\,\mathrm{d}t\big)\,g(0)$
as formal solution. This is the continuous analogue of accumulating a transport
matrix edge by edge, the recipe a discrete network will implement.

\Cref{fig:gauge:horizontal-vertical} pictures the splitting of $T_u P$ and the
role of the projection $\mathrm{d}\pi_u$.

\begin{figure}[ht]
\centering
\begin{tikzpicture}[>=Stealth, font=\small]
% base
\begin{scope}[shift={(0,-2)}]
  \draw[thick, fill=gdlblue!18]
    (-3,0) .. controls (-1,0.5) and (1,-0.3) .. (3,0.2) --
    (3,-0.5) .. controls (1,-1) and (-1,0) .. (-3,-0.5) -- cycle;
  \node[below] at (0,-0.8) {\textbf{base $\Mani$}};
  \fill[gdlblue!70] (0,0.1) circle (3pt);
  \node[below, font=\small] at (0,-0.1) {$x$};
\end{scope}
% fibre
\draw[thick, gdlorange!70, fill=gdlorange!18] (0,0) ellipse (0.4cm and 2cm);
\node[gdlorange!70, right, font=\small] at (0.5,1.5) {fibre $\pi^{-1}(x)\cong G$};
\fill[gdlred!70] (0,0.5) circle (4pt);
\node[right, font=\small, gdlred!70] at (0.2,0.5) {$u$};
% tangent decomposition
\begin{scope}[shift={(0,0.5)}]
  \draw[thick, dashed, gdlgray!40] (-2,-1) -- (2,1) -- (2.5,2) -- (-1.5,0) -- cycle;
  \node[gdlgray!60, font=\small] at (2.5,0) {$T_u P$};
  \draw[->, very thick, gdlpurple!70] (0,0) -- (0,1.2) node[above, font=\small] {$V_u P$};
  \node[gdlpurple!70, font=\tiny, left] at (-0.2,0.6) {vertical};
  \draw[->, very thick, gdlgreen!60!black] (0,0) -- (1.5,0.5) node[right, font=\small] {$H_u P$};
  \node[gdlgreen!60!black, font=\tiny, below] at (0.8,0.1) {horizontal};
  \draw[thick, dashed, gdlblue!50] (0,0) -- (1.5,1.7);
  \node[gdlblue!70, font=\tiny] at (1.0,1.3) {$v = v^V + v^H$};
\end{scope}
\draw[->, thick, gdlgray!70] (1.5,0.8) -- node[right, font=\small] {$\mathrm{d}\pi_u$} (1.5,-1.5);
\node[gdlgray!70, font=\tiny] at (2.7,-0.5) {$\ker(\mathrm{d}\pi_u)=V_u P$};
\node[font=\small, text width=5.2cm] at (6.6,0.2) {
  \textbf{Decomposition:}\\ $T_u P = V_u P \oplus H_u P$\\[0.4em]
  \textbf{Vertical:} $V_u P = \ker(\mathrm{d}\pi_u)$, tangent to the fibre
  (``gauge'' directions)\\[0.4em]
  \textbf{Horizontal:} $H_u P$ fixed by the connection (``spatial''
  directions)\\[0.4em]
  \textbf{$G$-invariance:} $H_{u\cdot g} = (R_g)_* H_u$
};
\end{tikzpicture}
\caption[Horizontal and vertical decomposition]{The connection splits each
tangent space $T_u P$ into a vertical part $V_u P = \ker(\mathrm{d}\pi_u)$,
tangent to the fibre, and a horizontal complement $H_u P$. The projection
$\mathrm{d}\pi_u$ kills the vertical directions; the horizontal directions
project isomorphically to the base.}
\label{fig:gauge:horizontal-vertical}
\end{figure}

\section{Gauge symmetry, fields, and equivariance}
\label{sec:gauge:symmetry}

We can now make the central symmetry precise. A choice of gauge is a choice of
frame at every point; changing it is a gauge transformation.

\begin{definition}[Gauge transformation]
\label{def:gauge:transformation}
A \emph{gauge transformation} is a smooth map $g\colon \Mani \to G$ assigning to
each point $p$ an element $g(p)$ of the structure group. The set of all such
maps forms the \emph{gauge group}
\[
  \mathcal{G} = C^\infty(\Mani, G)
\]
under pointwise multiplication, $(g_1 \cdot g_2)(p) = g_1(p)\, g_2(p)$.
\end{definition}

A gauge transformation rotates the local frame independently at every point.
\Cref{fig:gauge:freedom} shows two gauges on the same surface: the same
geometric data, described by different numerical components.

\begin{figure}[ht]
\centering
\begin{tikzpicture}[>=Stealth]
\def\surfacepath{(-2.4,-1.2) .. controls (-2.7,0.5) and (-1.9,2.2) .. (0,2.4)
  .. controls (1.9,2.6) and (2.9,1.2) .. (2.7,-0.2)
  .. controls (2.5,-1.5) and (0.7,-2.0) .. (-0.5,-1.8)
  .. controls (-1.6,-1.6) and (-2.1,-1.7) .. (-2.4,-1.2)}
\newcommand{\flen}{0.5}
% Gauge 1
\begin{scope}[shift={(-4.6,0)}]
  \node[above, font=\bfseries, text=gdlblue!70!black] at (0,2.9) {Gauge 1};
  \fill[gdlblue!15] \surfacepath; \draw[thick, gdlblue!40] \surfacepath;
  \foreach \p in {(-1.2,1.0),(1.0,1.2),(-0.9,-0.6),(1.1,-0.4),(0.0,0.3)} {
    \draw[->, thick, gdlgreen!50!black] \p -- ++(0:\flen);
    \draw[->, thick, gdlgreen!50!black] \p -- ++(90:\flen);
    \fill[gdlred!70] \p circle (2pt);
  }
\end{scope}
% Gauge 2
\begin{scope}[shift={(4.6,0)}]
  \node[above, font=\bfseries, text=gdlorange!70!black] at (0,2.9) {Gauge 2};
  \fill[gdlblue!15] \surfacepath; \draw[thick, gdlblue!40] \surfacepath;
  \draw[->, thick, gdlorange!70] (-1.2,1.0) -- ++(35:\flen);
  \draw[->, thick, gdlorange!70] (-1.2,1.0) -- ++(125:\flen);
  \draw[->, thick, gdlorange!70] (1.0,1.2) -- ++(200:\flen);
  \draw[->, thick, gdlorange!70] (1.0,1.2) -- ++(290:\flen);
  \draw[->, thick, gdlorange!70] (-0.9,-0.6) -- ++(110:\flen);
  \draw[->, thick, gdlorange!70] (-0.9,-0.6) -- ++(200:\flen);
  \draw[->, thick, gdlorange!70] (1.1,-0.4) -- ++(250:\flen);
  \draw[->, thick, gdlorange!70] (1.1,-0.4) -- ++(340:\flen);
  \draw[->, thick, gdlorange!70] (0.0,0.3) -- ++(60:\flen);
  \draw[->, thick, gdlorange!70] (0.0,0.3) -- ++(150:\flen);
  \foreach \p in {(-1.2,1.0),(1.0,1.2),(-0.9,-0.6),(1.1,-0.4),(0.0,0.3)}
    \fill[gdlred!70] \p circle (2pt);
\end{scope}
% arrow between
\draw[<->, very thick, gdlpurple!70] (-1.5,0) -- node[above, font=\small] {$g(x)\in\SO(2)$} (1.5,0);
\end{tikzpicture}
\caption[Gauge freedom]{Gauge freedom on a surface. The same surface admits
different choices of local frame (Gauge 1, with consistent orientation, and
Gauge 2, rotated independently at each point). A gauge-equivariant convolution
produces the same geometric result for either, the two being related by a
pointwise transformation $g(x) \in \SO(2)$.}
\label{fig:gauge:freedom}
\end{figure}

Features on the manifold are not bare numbers but components relative to the
chosen frame, so they must transform when the frame does. This is captured by
the type of a feature field, fixed by a representation $\rho$ of the structure
group.

\begin{definition}[Gauge-equivariant feature field]
\label{def:gauge:field}
Let $\rho\colon G \to \GL(V)$ be a representation. A field $f$ of \emph{type}
$\rho$ assigns to each point components in $V$ that, under a gauge transformation
$g$, transform as
\[
  f(p) \longmapsto \rho\big(g(p)^{-1}\big)\, f(p).
\]
Equivalently, regarding $f$ as a function on the total space, $f(u \cdot g) =
\rho(g^{-1})\, f(u)$ for all $u \in P$ and $g \in G$.
\end{definition}

Important special cases on a surface, with $G = \SO(2)$, are \emph{scalar}
fields (trivial $\rho$, invariant under gauge), \emph{vector} fields (the
standard $2$-dimensional representation, components rotating with the frame),
and complex fields transforming by $e^{i k\theta}$ for an integer ``spin'' $k$.
For a vector field with components $\mathbf{u} = u_1 e_1^v + u_2 e_2^v$, a frame
rotation by $\theta_v$ acts as
\[
  \begin{pmatrix} u_1' \\ u_2' \end{pmatrix} =
  \begin{pmatrix} \cos\theta_v & \sin\theta_v \\ -\sin\theta_v & \cos\theta_v
  \end{pmatrix}
  \begin{pmatrix} u_1 \\ u_2 \end{pmatrix},
\]
the standard transformation law of a type-$1$ field.

The design goal follows immediately. A network layer is \emph{gauge equivariant}
if it commutes with the gauge action: transforming the input frames produces the
correspondingly transformed output. If the final readout is a global scalar
prediction, the network is \emph{gauge invariant},
\[
  \Phi(g \cdot f) = \Phi(f) \qquad \text{for every gauge transformation } g.
\]
This guarantees that predictions are independent of the local parametrization or
chart, that vector-valued outputs rotate coherently with the frame, and --- in
the special case of node-permutation gauge on graphs --- that outputs are
invariant under relabelling of nodes. \Cref{fig:gauge:transformation} depicts a
gauge transformation acting on a graph signal. The connection to physics is
exact: gauge theories describe the fundamental forces through connections on
principal bundles with structure groups $\Ut(1)$, $\SU(2)$, $\SU(3)$; in
geometric deep learning the structure group encodes the local symmetries relevant
to the learning problem.

\begin{figure}[ht]
\centering
\begin{tikzpicture}[>=Stealth,
  vertex/.style={circle, draw, thick, minimum size=22pt}]
% original
\begin{scope}[shift={(-4.6,0)}]
  \node[above, font=\bfseries] at (1,2.4) {Features on a graph};
  \node[vertex, fill=gdlblue!20] (a) at (0,0) {$v_1$};
  \node[vertex, fill=gdlblue!20] (b) at (2,0) {$v_2$};
  \node[vertex, fill=gdlblue!20] (c) at (1,1.5) {$v_3$};
  \draw[thick] (a)--(b); \draw[thick] (b)--(c); \draw[thick] (c)--(a);
  \draw[->, thick, gdlred!70] (a) -- ++(-0.5,0.3) node[left, font=\tiny] {$f_1$};
  \draw[->, thick, gdlgreen!60!black] (b) -- ++(0.5,0.2) node[right, font=\tiny] {$f_2$};
  \draw[->, thick, gdlorange!70] (c) -- ++(0.2,0.5) node[above, font=\tiny] {$f_3$};
\end{scope}
\draw[->, very thick, gdlpurple!70] (-2,0.6) -- node[above, font=\small] {gauge $g:V\to G$} (2,0.6);
% transformed
\begin{scope}[shift={(4.6,0)}]
  \node[above, font=\bfseries] at (1,2.4) {Transformed features};
  \node[vertex, fill=gdlblue!20] (a2) at (0,0) {$v_1$};
  \node[vertex, fill=gdlblue!20] (b2) at (2,0) {$v_2$};
  \node[vertex, fill=gdlblue!20] (c2) at (1,1.5) {$v_3$};
  \draw[thick] (a2)--(b2); \draw[thick] (b2)--(c2); \draw[thick] (c2)--(a2);
  \draw[->, thick, gdlred!70] (a2) -- ++(0.3,0.5) node[left, font=\tiny] {$\rho(g_1)f_1$};
  \draw[->, thick, gdlgreen!60!black] (b2) -- ++(-0.2,0.5) node[right, font=\tiny] {$\rho(g_2)f_2$};
  \draw[->, thick, gdlorange!70] (c2) -- ++(-0.4,0.3) node[above, font=\tiny] {$\rho(g_3)f_3$};
  \draw[dashed, gdlpurple!50] (a2) ++(0.3,0) arc (0:60:0.3cm);
  \draw[dashed, gdlpurple!50] (b2) ++(0.3,0) arc (0:120:0.3cm);
  \draw[dashed, gdlpurple!50] (c2) ++(0.3,0) arc (0:150:0.3cm);
\end{scope}
\end{tikzpicture}
\caption[Gauge transformation on a graph]{A gauge transformation $g\colon V \to G$
rotates the reference frame independently at each vertex, sending the feature
$f_i$ to $\rho(g_i)\, f_i$. The dashed arcs indicate the local rotation applied at
each node.}
\label{fig:gauge:transformation}
\end{figure}

\begin{example}[Why gauge equivariance matters]
\label{ex:gauge:mesh}
On a $3$D triangle mesh $\Mani = (V, F)$ the tangent plane $T_v\Mani$ at each
vertex $v$ is well defined, but there is no canonical orthonormal basis
$(e_1^v, e_2^v)$ for it: any rotation $R_\theta \in \SO(2)$ gives an equally
valid frame. A naive multilayer perceptron predicting, say, the direction of
maximum curvature from the components $(u_1, u_2)$ of vectors in arbitrary local
coordinates fails silently. Suppose with the original frame it predicts the
direction $(\hat d_1, \hat d_2) = (0.8, 0.6)$; after rotating that frame by
$45^\circ$ --- a gauge transformation --- it might predict $(0.3, 0.9)$, an
entirely different direction. A gauge-equivariant network such as the Gauge
Equivariant Mesh CNN of \citet{cohen2019gaugeequivariantconvolutional} instead
guarantees $\hat{\mathbf{d}}' = R_\theta\,\hat{\mathbf{d}}$ under a frame rotation
$R_\theta$: the numerical components change, but the geometric direction
predicted does not.
\end{example}

\begin{example}[Tensorial fields on $3$D molecules]
\label{ex:gauge:molecules}
In modelling a molecule with atoms at positions $\{\mathbf{r}_i\}_{i=1}^n \subset
\mathbb{R}^3$ one often predicts tensorial quantities, such as the polarizability
tensor $\boldsymbol{\alpha} \in \mathbb{R}^{3\times 3}$. Under a global rotation
$R \in \SO(3)$ this tensor transforms as $\boldsymbol{\alpha}' = R\,
\boldsymbol{\alpha}\, R^{\!\top}$, a rank-$2$ tensor representation. Equivariant
networks such as NequIP~\citep{batzner2022nequip} represent intermediate features
as \emph{spherical tensors} transforming under irreducible representations of
$\SO(3)$ and build outputs through equivariant tensor products,
\[
  \boldsymbol{\alpha} = \sum_{\ell = 0, 2} W_\ell \big(\mathbf{Y}^{(\ell)} \otimes
  \mathbf{Y}^{(\ell)}\big),
\]
with $\mathbf{Y}^{(\ell)}$ spherical harmonics of order $\ell$ and $W_\ell$
learned weights. This guarantees the correct rotational behaviour automatically,
removing the need for rotational data augmentation and yielding physically
consistent predictions independent of the molecule's arbitrary orientation.
\end{example}

\section{Curvature, holonomy, and topological obstructions}
\label{sec:gauge:curvature}

Parallel transport along a path lets us compare frames at the endpoints, but the
result may depend on the path taken. The \emph{curvature} of a connection
measures this path dependence, equivalently the failure of the horizontal
distribution to be integrable.

\begin{definition}[Curvature and the Cartan structure equation]
\label{def:gauge:curvature}
The \emph{curvature} of a connection $\omega$ is the $\Lie{g}$-valued $2$-form
\[
  \Omega = \mathrm{d}\omega + \tfrac{1}{2}[\omega \wedge \omega],
\]
where $[\omega \wedge \omega](X, Y) = [\omega(X), \omega(Y)]$ uses the Lie
bracket; explicitly $\Omega(X, Y) = \mathrm{d}\omega(X, Y) + [\omega(X),
\omega(Y)]$. It is a horizontal, equivariant $2$-form. Rearranged, the definition
is the \emph{Cartan structure equation}
\[
  \mathrm{d}\omega = \Omega - \tfrac{1}{2}[\omega \wedge \omega].
\]
\end{definition}

\begin{theorem}[Bianchi identity]
\label{thm:gauge:bianchi}
The curvature satisfies $\mathrm{d}\Omega = [\omega \wedge \Omega]$.
\end{theorem}

\begin{proof}
Applying $\mathrm{d}$ to the definition of $\Omega$ and using $\mathrm{d}^2 = 0$,
\[
  \mathrm{d}\Omega = \mathrm{d}(\mathrm{d}\omega) + \tfrac{1}{2}\mathrm{d}[\omega
  \wedge \omega] = \tfrac{1}{2}\big([\mathrm{d}\omega \wedge \omega] - [\omega
  \wedge \mathrm{d}\omega]\big).
\]
Substituting $\mathrm{d}\omega = \Omega - \tfrac{1}{2}[\omega \wedge \omega]$ and
using the Jacobi identity for the Lie bracket, the quadratic terms cancel,
leaving $\mathrm{d}\Omega = [\omega \wedge \Omega]$.
\end{proof}

\begin{definition}[Holonomy]
\label{def:gauge:holonomy}
For a loop $\gamma$ based at $p$ (so $\gamma(0) = \gamma(1) = p$), the
\emph{holonomy} is the element $g_\gamma \in G$ with $\Gamma_\gamma(u) = u \cdot
g_\gamma$ for all $u \in \pi^{-1}(p)$. The \emph{holonomy group}
$\Hol_p(\omega)$ is the subgroup of $G$ generated by the holonomies of all loops
based at $p$.
\end{definition}

Curvature is holonomy in the infinitesimal limit: for a small loop enclosing
area $\mathrm{d}A$ the holonomy is approximately $\exp(\Omega \cdot \mathrm{d}A)$,
and the Ambrose--Singer theorem makes this precise by identifying $\Hol_p(\omega)$
with the subgroup generated by the curvature values $\Omega(X, Y)$. When the
curvature vanishes the connection is \emph{flat} and parallel transport depends
only on the homotopy class of the path; a vector carried around a contractible
loop returns unchanged. On a curved domain the curvature is generally nonzero,
and this is exactly why a frame transported around a closed loop comes back
rotated --- the obstruction a gauge-equivariant convolution must respect.
\Cref{fig:gauge:holonomy} contrasts the flat and curved cases.

\begin{figure}[ht]
\centering
\begin{tikzpicture}[>=Stealth, font=\small]
% flat
\begin{scope}[shift={(-4.2,0)}]
  \node[above, font=\bfseries] at (0,2.6) {Flat connection ($\Omega=0$)};
  \fill[gdlblue!18] (-1.6,-1.6) rectangle (1.6,1.6);
  \draw[thick, gdlblue!30] (-1.6,-1.6) rectangle (1.6,1.6);
  \draw[very thick, gdlred!70] (0,0) -- (1.2,0) -- (1.2,1.2) -- (0,1.2) -- cycle;
  \draw[->, thick, gdlgreen!60!black] (0,0) -- (0.5,0.25);
  \draw[->, thick, gdlgreen!60!black] (1.2,0) -- (1.7,0.25);
  \draw[->, thick, gdlgreen!60!black] (1.2,1.2) -- (1.7,1.45);
  \draw[->, thick, gdlgreen!60!black] (0,1.2) -- (0.5,1.45);
  \node[font=\footnotesize] at (0,-2.0) {$v_{\text{final}} = v_{\text{initial}}$};
\end{scope}
% curved
\begin{scope}[shift={(4.2,0)}]
  \node[above, font=\bfseries] at (0,2.6) {Curved connection ($\Omega\neq0$)};
  \shade[ball color=gdlpurple!20, opacity=0.8] (0,0) circle (1.6cm);
  \draw[thick, gdlpurple!40] (0,0) circle (1.6cm);
  \draw[very thick, gdlred!70] (0,1.6) arc (90:0:1.6cm) arc (0:-90:1.6cm and 0.4cm) arc (-90:90:0.4cm and 1.6cm);
  \draw[->, thick, gdlgreen!60!black] (0,1.6) -- (0.45,1.85);
  \draw[->, thick, gdlgreen!60!black] (1.6,0) -- (1.9,0.45);
  \draw[->, thick, gdlgreen!60!black] (0,-0.4) -- (0.45,-0.15);
  \draw[->, very thick, gdlorange!70, dashed] (0,1.6) -- (-0.3,1.95);
  \draw[thick, gdlorange!70] (0.25,1.7) arc (30:120:0.3cm);
  \node[gdlorange!70, font=\tiny] at (-0.1,2.05) {$\theta$};
  \node[font=\footnotesize] at (0,-2.0) {$v_{\text{final}} \neq v_{\text{initial}}$};
\end{scope}
\node[font=\small, text width=8cm, align=center] at (0,-3.0)
  {\textbf{Ambrose--Singer:}\ $\Hol_p(\omega) = \big\langle \Omega(X,Y) : X,Y \in
  T_u P\big\rangle$ --- holonomy is generated by the curvature.};
\end{tikzpicture}
\caption[Holonomy and curvature]{Parallel transport around a closed loop. For a
flat connection (left) the transported vector returns unchanged and the holonomy
is trivial. For a curved connection (right) it returns rotated by an angle
$\theta$; the holonomy equals the integral of the curvature over the enclosed
surface.}
\label{fig:gauge:holonomy}
\end{figure}

Curvature also explains \emph{why} gauge equivariance is necessary rather than
merely convenient.

\begin{proposition}[Topological obstruction to a global gauge]
\label{prop:gauge:obstruction}
A principal bundle $P \to \Mani$ admits a global section --- a consistent choice
of gauge over all of $\Mani$ --- if and only if it is trivial,
$P \cong \Mani \times G$. In particular:
\begin{enumerate}
  \item the orthonormal frame bundle of $S^2$ admits no global section (by the
    hairy-ball theorem);
  \item every principal bundle over a contractible manifold is trivial.
\end{enumerate}
\end{proposition}

\begin{proof}
A global section $s\colon \Mani \to P$ induces the trivialization
$\Mani \times G \to P$, $(p, g) \mapsto s(p) \cdot g$, an isomorphism of bundles
by the free and transitive action of $G$ on each fibre; conversely a
trivialization $P \cong \Mani \times G$ provides the section $s(p) = (p, e)$.
For (1), a section of the frame bundle would supply a continuous field of
orthonormal frames, in particular a nowhere-vanishing tangent vector field,
contradicting the hairy-ball theorem ($\chi(S^2) = 2 \neq 0$). For (2), a
contractible base is homotopy-equivalent to a point, and principal bundles are
classified by homotopy classes of maps into the classifying space $BG$, so every
bundle over it is trivial.
\end{proof}

The consequence for learning is decisive: on $S^2$ and most manifolds of
interest there is no globally consistent coordinate system, and any global
parametrization has singularities (the poles in spherical coordinates). Gauge
equivariance frees the network from these arbitrary local choices and avoids
artefacts at the singularities, instead of papering over them with a fixed
chart.

\section{Gauge-equivariant convolution}
\label{sec:gauge:conv}

We now assemble the pieces into a convolution. The difficulty a gauge-equivariant
layer must overcome is concrete: to combine the feature at a point $p$ with the
features of its neighbours, those neighbours' features live in different fibres,
expressed in different local frames, and cannot be added directly. The
connection resolves this --- parallel transport carries each neighbour's frame
to $p$ before the features are combined. We first formalize the distinction
between global and local transformations underlying the construction.

\begin{definition}[Local versus global transformation]
\label{def:gauge:local-global}
For a group $G$ and a field $f\colon \Mani \to \mathbb{R}^d$:
a \emph{global} transformation is $T_g[f](x) = \rho(g)\, f(g^{-1} \cdot x)$, with
a single $g \in G$ used at every point $x$; a \emph{local} (gauge)
transformation is $T_{g(\cdot)}[f](x) = \rho(g(x))\, f(\phi_{g(x)}(x))$, with
$g\colon \Mani \to G$ assigning a possibly different element to each point.
\end{definition}

\begin{theorem}[Gauge-equivariant convolution]
\label{thm:gauge:conv}
Let $f$ be a feature field on $\Mani$ and let $\Gamma_{q \to p}$ denote parallel
transport from $q$ to $p$. The convolution
\[
  (f \star \psi)(p) = \int_{\mathcal{N}(p)}
    \psi\big(q, p\big)\; \Gamma_{q \to p}\, f(q)\; \mathrm{d}q,
\]
which transports each neighbour's feature into the frame at $p$ before applying
the learnable kernel $\psi$, maps gauge-equivariant fields to
gauge-equivariant fields.
\end{theorem}

\begin{proof}
Let $\rho_p$ denote a gauge change at $p$, so that feature fields transform as
$f(q) \mapsto \rho_q\, f(q)$. The transport operator $\Gamma_{q \to p}$ changes
basis from the frame at $q$ to the frame at $p$, hence transforms as
$\Gamma_{q \to p} \mapsto \rho_p\, \Gamma_{q \to p}\, \rho_q^{-1}$. Substituting,
\[
  (f \star \psi)(p) \mapsto \int_{\mathcal{N}(p)} \psi(q,p)\,
    \big(\rho_p\, \Gamma_{q\to p}\, \rho_q^{-1}\big)\,\big(\rho_q\, f(q)\big)\,
    \mathrm{d}q
  = \rho_p \int_{\mathcal{N}(p)} \psi(q,p)\, \Gamma_{q\to p}\, f(q)\, \mathrm{d}q,
\]
since $\rho_q^{-1}\rho_q = \mathrm{Id}$. Thus $(f \star \psi)(p) \mapsto \rho_p\,
(f \star \psi)(p)$, the transformation law of an equivariant
field~\citep{cohen2019gaugeequivariantconvolutional}.
\end{proof}

The cancellation is the same algebraic mechanism that makes the field strength
$F_{\mu\nu} = \partial_\mu A_\nu - \partial_\nu A_\mu$ gauge invariant in
electromagnetism even though the potential $A_\mu$ is not: gauge factors at the
source point cancel against the inverse factors carried by the connection.
Crucially, gauge equivariance imposes \emph{linear} constraints on the kernel
$\psi$, cutting the parameter space down to an affine subspace that can be solved
for analytically before training --- the same kernel-constraint principle that
makes steerable group convolutions practical in \Cref{ch:groups}.

\subsection{The icosahedral CNN}
\label{sec:gauge:icosahedral}

The gap between the smooth theory and a finite, trainable network was first
bridged in practice by \citet{cohen2019gaugeequivariantconvolutional}, who chose
the icosahedron as a discretization of the sphere $S^2$. The choice is far from
arbitrary: among the regular polyhedra the icosahedron best approximates the
sphere, it has a rich symmetry group with $120$ elements, its vertices are
nearly uniformly spaced, and at only $12$ vertices, $20$ faces, and $30$ edges
(refinable by geodesic subdivision) it is computationally light.

On this mesh the structure group is $\SO(2)$ --- planar rotations of the tangent
frame at each vertex. A discrete feature field $f\colon V \to V_\rho$ lives on
the vertices, and the gauge-equivariant convolution becomes a local sum over
neighbours,
\[
  (f \star \psi)(v_i) = \sum_{j \in \mathcal{N}(i)}
    \psi_{ij}\; \Gamma_{ij}\, f(v_j),
\]
where $\Gamma_{ij} \in \SO(2)$ is the discrete parallel transport from vertex
$v_j$ to $v_i$ along the connecting edge --- a rotation by a fixed angle fixed by
the mesh geometry --- and $\psi_{ij}$ are the learnable kernel weights,
constrained so that the layer respects the $\SO(2)$ gauge action. The transport
of a feature from a distant vertex is obtained by composing the per-edge
rotations $\Gamma_{ij} = \exp(-\omega_e)$ along a geodesic path, exactly the
discrete counterpart of the path-ordered exponential of
\Cref{sec:gauge:transport}: initialize the transport to the identity and multiply
by $\exp(-\omega_e)$ for each edge $e$ traversed. To avoid degenerate solutions
one penalizes excess curvature through a regularizer summing the edge angles
around each triangle,
\[
  \mathcal{R}(\theta) = \sum_{\text{triangles } T}
    \Big|\sum_{(i,j) \in \partial T} \theta_{ij}\Big|^2,
\]
a discrete reading of the Bianchi identity of \Cref{thm:gauge:bianchi}: it
controls the holonomy accumulated around small loops.

\Cref{fig:gauge:icosahedral} sketches the mechanism. The decisive practical
benefit is that this whole computation can be expressed in terms of standard
convolution-like operations on a flattened representation of the icosahedron,
making the network as efficient as an ordinary CNN while being \emph{exactly}
gauge equivariant --- and, because the icosahedral symmetry group acts on the
mesh, equivariant to that discrete subgroup of sphere rotations as well.

\begin{figure}[ht]
\centering
\begin{tikzpicture}[>=Stealth, font=\small]
% central vertex and neighbours (hexagonal local patch)
\coordinate (c) at (0,0);
\foreach \a/\name in {0/n0,60/n1,120/n2,180/n3,240/n4,300/n5}
  \coordinate (\name) at (\a:1.8);
% edges
\foreach \name in {n0,n1,n2,n3,n4,n5} \draw[gdlgray!55, thick] (c) -- (\name);
\foreach \a/\b in {n0/n1,n1/n2,n2/n3,n3/n4,n4/n5,n5/n0}
  \draw[gdlgray!35] (\a) -- (\b);
% transported frames (rotated) at neighbours
\foreach \name/\rot in {n0/20,n1/-15,n2/40,n3/10,n4/-30,n5/25} {
  \draw[->, thick, gdlorange!75] (\name) -- ++(\rot:0.6);
  \draw[->, thick, gdlorange!75] (\name) -- ++({\rot+90}:0.6);
  \fill[gdlblue!65] (\name) circle (2.5pt);
}
% transport arrows toward centre
\foreach \name in {n0,n2,n4}
  \draw[->, gdlpurple!70, thick, dashed] ($(\name)!0.30!(c)$) -- ($(\name)!0.62!(c)$);
% central frame
\draw[->, very thick, gdlgreen!55!black] (c) -- ++(0:0.7);
\draw[->, very thick, gdlgreen!55!black] (c) -- ++(90:0.7);
\fill[gdlred!75] (c) circle (3pt);
\node[below right, font=\footnotesize, text=gdlgreen!45!black] at (c) {$v_i$};
\node[font=\footnotesize, text=gdlpurple!70] at (1.15,0.55) {$\Gamma_{ij}$};
\node[font=\footnotesize, text=gdlorange!70!black] at (-1.9,1.55) {$f(v_j)$};
\end{tikzpicture}
\caption[Gauge-equivariant convolution on a mesh]{Local gauge-equivariant
convolution at a vertex $v_i$. Each neighbour carries a feature in its own local
frame (orange). Before combining, the connection parallel-transports it to the
frame at $v_i$ through the rotation $\Gamma_{ij} \in \SO(2)$ (purple), after
which the learnable kernel $\psi_{ij}$ is applied. The result transforms
coherently under any change of local frame.}
\label{fig:gauge:icosahedral}
\end{figure}

\subsection{The general framework and outlook}
\label{sec:gauge:outlook}

The icosahedral case instantiates a general recipe. Choosing a base manifold
$\Mani$ for the data and a structure group $G$ for the relevant local symmetries
fixes a principal bundle; a connection supplies parallel transport; and the
gauge-equivariant convolution of \Cref{thm:gauge:conv} produces a layer whose
output type is governed by a chosen representation $\rho$. Different problems
call for different bundles: trivial bundles when symmetries are globally
consistent, frame bundles when invariance under local coordinate changes is
required, spin bundles for data with orientation structure, and richer gauge
bundles for fields with internal symmetries (electromagnetism with $\Ut(1)$,
Yang--Mills theories with non-abelian groups, lattice quantum chromodynamics with
$\SU(3)$). Under mild conditions such networks are universal approximators of
gauge-equivariant maps: the space of realizable functions is dense in the
$L^2$ space of gauge-equivariant fields, with approximation rate
$O(n^{-2/(d_\Mani + d_G)})$ governed by the dimensions of the base and the group
and constants depending on the curvature of the bundle.

This framework is the capstone of the five Gs. It strictly contains the others:
a global symmetry group is the special case of a trivial bundle with a flat
connection, recovering the group-equivariant networks of \Cref{ch:groups}; a
manifold with its Levi-Civita connection recovers the geodesic convolutions of
\Cref{ch:geodesics} as a particular gauge choice; and the spherical networks
that motivated this development sit inside it as the icosahedral instance. The
nesting
\[
  \mathcal{F}_{\text{Euclidean}} \subset \mathcal{F}_{\text{group}}
  \subset \mathcal{F}_{\text{geodesic}} \subset \mathcal{F}_{\text{gauge}}
\]
summarizes how each of the previous four Gs is a constrained version of
gauge-equivariant convolution.

\section{Clifford and geometric algebras}
\label{sec:gauge:clifford}

The gauge viewpoint admits a powerful algebraic generalization. Where a feature
field expresses a directional quantity in components relative to a frame, a
\emph{geometric (Clifford) algebra} lets scalars, vectors, oriented areas, and
oriented volumes live in a single graded object and transform together under
rotations and reflections. This unifies complex numbers, quaternions, and tensor
representations, and provides equivariant building blocks for both convolution
and attention.

\begin{definition}[Clifford algebra]
\label{def:gauge:clifford}
Let $V$ be a finite-dimensional real vector space with quadratic form $Q$. The
\emph{Clifford algebra} $\Cl(V, Q)$ is the unital associative algebra generated
by $V$ subject to $v^2 = Q(v)\, 1$ for all $v \in V$. For a metric of signature
$(p, q)$ we write $\Cl(p, q)$, with $p$ positive-norm and $q$ negative-norm
directions.
\end{definition}

The algebra unifies several familiar structures: $\Cl(0,0) \cong \mathbb{R}$,
$\Cl(0,1) \cong \mathbb{C}$, $\Cl(0,2) \cong \mathbb{H}$ (the quaternions), and
$\Cl(3,0)$ is the Euclidean geometric algebra of $\mathbb{R}^3$.

\begin{definition}[Multivector and grading]
\label{def:gauge:multivector}
Elements of $\Cl(p, q)$ are \emph{multivectors} and decompose uniquely by grade.
With $n = p + q$,
\[
  M = \langle M \rangle_0 + \langle M \rangle_1 + \cdots + \langle M \rangle_n,
\]
where $\langle M \rangle_0 \in \mathbb{R}$ is a scalar, $\langle M \rangle_1 \in
\operatorname{span}\{\mathbf{e}_1, \ldots, \mathbf{e}_n\}$ a vector,
$\langle M \rangle_2$ a bivector (oriented area), $\langle M \rangle_k$ a
$k$-vector of dimension $\binom{n}{k}$, and $\langle M \rangle_n$ the
pseudoscalar.
\end{definition}

\begin{theorem}[Dimension of the algebra]
\label{thm:gauge:clifford-dim}
The Clifford algebra $\Cl(p, q)$ is a vector space of dimension $2^{p+q}$, with
canonical basis the ordered products of subsets of
$\{\mathbf{e}_1, \ldots, \mathbf{e}_{p+q}\}$:
$\dim \Cl(p,q) = \sum_{k=0}^{p+q} \binom{p+q}{k} = 2^{p+q}$.
\end{theorem}

\begin{proof}
Let $n = p+q$ and $\{\mathbf{e}_i\}$ an orthonormal basis. The basis multivectors
are the ordered products $\mathbf{e}_{i_1}\cdots\mathbf{e}_{i_k}$ with
$1 \le i_1 < \cdots < i_k \le n$, $0 \le k \le n$ (with $k=0$ the scalar $1$).
These are linearly independent: the relations $\mathbf{e}_i\mathbf{e}_j =
-\mathbf{e}_j\mathbf{e}_i$ for $i \neq j$ order any product, and
$\mathbf{e}_i^2 = \pm 1$ removes squares. The number of subsets of $\{1, \ldots,
n\}$ is $\sum_{k=0}^{n}\binom{n}{k} = 2^n$, the dimension of the algebra.
\end{proof}

\begin{definition}[Geometric product]
\label{def:gauge:geometric-product}
For vectors $\mathbf{u}, \mathbf{v} \in V$ the \emph{geometric product} unifies
the inner and exterior products:
\[
  \mathbf{u}\mathbf{v} = \mathbf{u}\cdot\mathbf{v} + \mathbf{u}\wedge\mathbf{v},
  \qquad
  \mathbf{u}\cdot\mathbf{v} = \tfrac{1}{2}(\mathbf{u}\mathbf{v} +
  \mathbf{v}\mathbf{u}),
  \quad
  \mathbf{u}\wedge\mathbf{v} = \tfrac{1}{2}(\mathbf{u}\mathbf{v} -
  \mathbf{v}\mathbf{u}).
\]
\end{definition}

\begin{example}[Rotors and reflections in $\Cl(3,0)$]
\label{ex:gauge:rotors}
In $\Cl(2,0)$ a rotation by angle $\theta$ is encoded by the \emph{rotor}
$R = \cos(\theta/2) + \sin(\theta/2)\,\mathbf{e}_1\wedge\mathbf{e}_2 =
e^{(\theta/2)\mathbf{e}_1\wedge\mathbf{e}_2}$, acting on a vector by the sandwich
product $\mathbf{v}' = R\,\mathbf{v}\,\widetilde{R}$, where $\widetilde{R}$ is the
reverse. A reflection across the hyperplane orthogonal to a unit vector
$\mathbf{n}$ is $\mathbf{v}' = -\mathbf{n}\,\mathbf{v}\,\mathbf{n}$; for instance
$\mathbf{n} = \mathbf{e}_1$ and $\mathbf{v} = a\mathbf{e}_1 + b\mathbf{e}_2$ give
$\mathbf{v}' = -a\mathbf{e}_1 + b\mathbf{e}_2$, flipping the perpendicular
component. A bivector squares to $-1$
($(\mathbf{e}_1\wedge\mathbf{e}_2)^2 = -1$), so it plays the algebraic role of the
imaginary unit. These sandwich operations are exactly gauge transformations
realized inside the algebra, with rotors implementing $\SO(n)$ and reflections
generating the full $\Ot(n)$.
\end{example}

\begin{definition}[Clifford convolution and Clifford bundle]
\label{def:gauge:clifford-conv}
For a Clifford-valued signal $f\colon \Omega \to \Cl(p,q)$ and a kernel
$\psi\colon G \to \Cl(p,q)$ on a transformation group $G$, the \emph{Clifford
convolution} is
\[
  (f \star \psi)(x) = \int_G \rho_g^{-1}\big(\psi(g)\big)\cdot f(g^{-1}\cdot x)\,
  \mathrm{d}g,
\]
with $\cdot$ the geometric product and $\rho_g$ the action of $G$ on the algebra.
On a Riemannian manifold $(\Mani, g)$ the \emph{Clifford bundle} $\Cl(T\Mani)$ has
fibre $\Cl(T_p\Mani, g_p)$ at $p$, and feature fields are its sections: $s(p) \in
\Cl(T_p\Mani)$.
\end{definition}

Clifford-valued convolutions package several geometric aspects into one operation:
the grades of a multivector carry distinct, interpretable geometric meaning;
equivariance is automatic from the algebraic structure; and the operations are
amenable to vectorized implementation. They subsume steerable group convolutions
(\Cref{ch:groups}): when a representation $\rho$ decomposes into irreducibles
isomorphic to fixed-grade multivector spaces of $\Cl(n,0)$, equivariant linear
maps correspond to grade-preserving Clifford operations, recovering the block
structure of steerable CNNs. The choice of algebra is itself a modelling
decision: the Euclidean $\Cl(3,0)$ ($8$ dimensions) is light but limited to
$\SO(3)$; the projective and conformal algebras $\Cl(4,1)$ ($32$ dimensions)
represent translations, spheres, and inversions at higher cost, and empirically
the conformal algebra performs best~\citep{dehaan2024euclidean}. Learning the
metric of the algebra adaptively (Clifford-group-equivariant networks with
learnable $G_\theta = \sum_i \lambda_i(\theta)\,\mathbf{v}_i\mathbf{v}_i^{\!\top}$)
lets the geometry of the algebra adjust to the data while preserving
equivariance~\citep{ali2024metric}.

\paragraph{Geometric Algebra Transformers (GATr).}
The same algebra powers equivariant attention. GATr~\citep{brehmer2023geometric}
operates in the projective geometric algebra $\Cl(3,0,1)$ (a $16$-dimensional
algebra over the basis $\{\mathbf{e}_0, \mathbf{e}_1, \mathbf{e}_2,
\mathbf{e}_3\}$ with $\mathbf{e}_i^2 = +1$ for $i = 1,2,3$ and $\mathbf{e}_0^2 =
0$). The degenerate direction $\mathbf{e}_0$ is what makes affine geometry ---
including translations --- expressible, not just linear geometry. Geometric
objects of $\mathbb{R}^3$ receive canonical multivector representations: scalars
at grade $0$, planes $\pi = n_1\mathbf{e}_1 + n_2\mathbf{e}_2 + n_3\mathbf{e}_3 +
d\,\mathbf{e}_0$ at grade $1$, lines at grade $2$, points at grade $3$, and the
pseudoscalar $I = \mathbf{e}_0\wedge\mathbf{e}_1\wedge\mathbf{e}_2\wedge
\mathbf{e}_3$ at grade $4$, with grade dimensions $1, 4, 6, 4, 1$ summing to
$16$. Rotations and translations act by rotors through the sandwich product, so
attention built from grade-preserving linear maps, geometric products, and an
equivariant inner product is exactly $\mathrm{E}(3)$-equivariant. GATr is the
attention-mechanism counterpart of the gauge-equivariant convolution: the gauge
group is realized algebraically, by Clifford operations, rather than by an
explicit connection.

\section{Topological extensions}
\label{sec:gauge:topology}

The bundle picture treats the domain as a smooth manifold, but many domains are
genuinely combinatorial, with cells of several dimensions and incidence relations
richer than the edges of a graph. \emph{Topological deep learning} extends
message passing to such structures, and the algebra of boundary operators plays
the role the connection played above: it specifies how information flows between
cells of adjacent dimension.

\subsection{Cell complexes, chains, and boundary operators}
\label{sec:gauge:chains}

Recall that a \emph{simplicial complex} $\mathcal{K}$ is a downward-closed family
of finite sets meeting pairwise in common faces; a $0$-simplex is a vertex, a
$1$-simplex an edge, a $2$-simplex a triangle, a $3$-simplex a tetrahedron. A
\emph{cell (CW) complex} generalizes this by gluing cells $e^n$ (homeomorphic to
open balls) along characteristic maps $\phi\colon \partial\mathbb{D}^n \to
X^{n-1}$ to the lower skeleton; cells may have arbitrary shape (a cube is a single
$3$-cell rather than several tetrahedra). A \emph{combinatorial complex} relaxes
further to any family of cells with a partial inclusion relation, subsuming
graphs, hypergraphs, simplicial and cell complexes~\citep{hajij2023topological}.

\begin{definition}[Chains and the boundary operator]
\label{def:gauge:chains}
The space of \emph{$k$-chains} is $C_k(\mathcal{K}) = \{\sum_i a_i\sigma_i : a_i
\in \mathbb{R}, \sigma_i\ \text{a } k\text{-simplex}\}$. The \emph{boundary
operator} $\partial_k\colon C_k \to C_{k-1}$ acts on a $k$-simplex $\sigma =
[v_0, \ldots, v_k]$ by
\[
  \partial_k \sigma = \sum_{i=0}^k (-1)^i [v_0, \ldots, \hat v_i, \ldots, v_k],
\]
with $\hat v_i$ denoting omission, extended linearly. The \emph{coboundary}
operator is the adjoint $\delta_k = \partial_{k+1}^{\!\top}\colon C_k \to
C_{k+1}$.
\end{definition}

\begin{theorem}[$\partial^2 = 0$]
\label{thm:gauge:partial-squared}
For all $k$, $\partial_{k-1}\circ\partial_k = 0$: the boundary of a boundary is
empty.
\end{theorem}

\begin{proof}
For $\sigma = [v_0, \ldots, v_k]$,
\[
  \partial_{k-1}(\partial_k\sigma) = \sum_{i} (-1)^i \Big(\sum_{j<i}(-1)^j
  [\ldots \hat v_j \ldots \hat v_i \ldots] + \sum_{j>i}(-1)^{j-1}
  [\ldots \hat v_i \ldots \hat v_j \ldots]\Big).
\]
Each face with two vertices omitted appears exactly twice, once from each order
of removal, with opposite signs, so all terms cancel.
\end{proof}

For example $\partial_2[v_0, v_1, v_2] = [v_1, v_2] - [v_0, v_2] + [v_0, v_1]$,
the three oriented edges of the triangle, and applying $\partial_1$ gives zero.
\Cref{fig:gauge:simplicial-hierarchy} shows the resulting hierarchy of chains
linked by boundary and coboundary maps.

\begin{figure}[ht]
\centering
\begin{tikzpicture}[>=Stealth, font=\small, scale=0.62, transform shape]
\def\leveltop{7}\def\levelmid{0}\def\levelbot{-7}
\begin{scope}[on background layer]
  \fill[gdlblue!6, rounded corners=8pt] (-2.6,\leveltop-2.2) rectangle (7.8,\leveltop+2.2);
  \fill[gdlorange!6, rounded corners=8pt] (-2.6,\levelmid-2.2) rectangle (7.8,\levelmid+2.2);
  \fill[gdlred!6, rounded corners=8pt] (-2.6,\levelbot-2.2) rectangle (7.8,\levelbot+2.2);
\end{scope}
% level 0: nodes
\begin{scope}[shift={(0,\leveltop)}]
  \node[font=\bfseries, text=gdlblue!70!black, anchor=west] at (-2.4,1.5) {$0$-chains (nodes), $L_0$};
  \foreach \a/\b in {0/0,1.8/0.6,3.2/-0.2,1.0/-1.5,2.8/-1.8,4.2/-1.2}
    \fill[gdlblue!70] (\a,\b) circle (4pt);
  \draw[gdlgray!40,thick] (0,0)--(1.8,0.6) (0,0)--(1.0,-1.5) (1.8,0.6)--(3.2,-0.2)
    (1.8,0.6)--(1.0,-1.5) (1.8,0.6)--(2.8,-1.8) (3.2,-0.2)--(2.8,-1.8)
    (3.2,-0.2)--(4.2,-1.2) (1.0,-1.5)--(2.8,-1.8) (2.8,-1.8)--(4.2,-1.2);
\end{scope}
% level 1: edges
\begin{scope}[shift={(0,\levelmid)}]
  \node[font=\bfseries, text=gdlorange!70!black, anchor=west] at (-2.4,1.5) {$1$-chains (edges), $L_1$};
  \draw[gdlorange!70,very thick] (0,0)--(1.8,0.6) (0,0)--(1.0,-1.5) (1.8,0.6)--(3.2,-0.2)
    (1.8,0.6)--(1.0,-1.5) (1.8,0.6)--(2.8,-1.8) (3.2,-0.2)--(2.8,-1.8)
    (3.2,-0.2)--(4.2,-1.2) (1.0,-1.5)--(2.8,-1.8) (2.8,-1.8)--(4.2,-1.2);
  \foreach \a/\b in {0/0,1.8/0.6,3.2/-0.2,1.0/-1.5,2.8/-1.8,4.2/-1.2}
    \fill[gdlgray!50] (\a,\b) circle (3pt);
\end{scope}
% level 2: faces
\begin{scope}[shift={(0,\levelbot)}]
  \node[font=\bfseries, text=gdlred!70!black, anchor=west] at (-2.4,1.5) {$2$-chains (faces), $L_2$};
  \fill[gdlred!15] (0,0)--(1.8,0.6)--(1.0,-1.5)--cycle;
  \fill[gdlred!15] (1.8,0.6)--(1.0,-1.5)--(2.8,-1.8)--cycle;
  \fill[gdlred!15] (3.2,-0.2)--(2.8,-1.8)--(4.2,-1.2)--cycle;
  \draw[gdlred!60,thick] (0,0)--(1.8,0.6)--(1.0,-1.5)--cycle
    (1.8,0.6)--(1.0,-1.5)--(2.8,-1.8)--cycle (3.2,-0.2)--(2.8,-1.8)--(4.2,-1.2)--cycle;
  \foreach \a/\b in {0/0,1.8/0.6,3.2/-0.2,1.0/-1.5,2.8/-1.8,4.2/-1.2}
    \fill[gdlgray!50] (\a,\b) circle (3pt);
\end{scope}
% boundary / coboundary arrows
\draw[->, gdlpurple!70, line width=1.5pt] (-3.4,\levelmid+2.3) -- (-3.4,\leveltop-2.3);
\node[gdlpurple!70!black] at (-3.4,{(\levelmid+\leveltop)/2}) {$\partial_1$};
\draw[->, gdlgreen!70!black, line width=1.5pt] (-5.4,\leveltop-2.3) -- (-5.4,\levelmid+2.3);
\node[gdlgreen!70!black] at (-5.4,{(\leveltop+\levelmid)/2}) {$\delta_0$};
\draw[->, gdlpurple!70, line width=1.5pt] (-3.4,\levelbot+2.3) -- (-3.4,\levelmid-2.3);
\node[gdlpurple!70!black] at (-3.4,{(\levelbot+\levelmid)/2}) {$\partial_2$};
\draw[->, gdlgreen!70!black, line width=1.5pt] (-5.4,\levelmid-2.3) -- (-5.4,\levelbot+2.3);
\node[gdlgreen!70!black] at (-5.4,{(\levelmid+\levelbot)/2}) {$\delta_1$};
\end{tikzpicture}
\caption[Simplicial hierarchy]{Multi-order message passing on a simplicial
complex: signals on nodes ($0$-chains), edges ($1$-chains), and faces
($2$-chains) are linked by boundary operators $\partial_k$ and coboundary
operators $\delta_k$. The Hodge Laplacian $L_k$ couples the level-$k$ signal to
its neighbours of dimension $k \pm 1$.}
\label{fig:gauge:simplicial-hierarchy}
\end{figure}

\subsection{Hodge Laplacians and homology}
\label{sec:gauge:hodge}

\begin{definition}[Hodge Laplacian]
\label{def:gauge:hodge-laplacian}
The \emph{Hodge Laplacian} of order $k$ is
\[
  L_k = \partial_{k+1}\partial_{k+1}^{\!\top} + \partial_k^{\!\top}\partial_k
      = L_k^{\text{up}} + L_k^{\text{down}},
\]
with $L_k^{\text{down}} = \partial_k^{\!\top}\partial_k$ the flow from
$(k-1)$-cells and $L_k^{\text{up}} = \partial_{k+1}\partial_{k+1}^{\!\top}$ the
flow to $(k+1)$-cells. For $k = 0$ it reduces to the graph Laplacian.
\end{definition}

\begin{theorem}[Properties of the Hodge Laplacian]
\label{thm:gauge:hodge-props}
$L_k$ is symmetric and positive semidefinite; its kernel is isomorphic to the
$k$-th homology, $\ker(L_k) \cong H_k(\mathcal{K})$; and $C_k$ decomposes
orthogonally as
\[
  C_k = \Img(\partial_{k+1}) \oplus \ker(L_k) \oplus \Img(\partial_k^{\!\top}).
\]
\end{theorem}

\begin{proof}
Symmetry and positivity are immediate, since $L_k$ is a sum of terms $A^{\!\top}
A$ and $\langle L_k\mathbf{x}, \mathbf{x}\rangle = \|\partial_k\mathbf{x}\|^2 +
\|\partial_{k+1}^{\!\top}\mathbf{x}\|^2 \ge 0$. Hence $L_k\mathbf{x} = 0$ iff
$\partial_k\mathbf{x} = 0$ and $\partial_{k+1}^{\!\top}\mathbf{x} = 0$, that is,
$\mathbf{x}$ is a $k$-cycle orthogonal to $\Img(\partial_{k+1})$; this is exactly
the space of harmonic representatives of $H_k = \ker\partial_k /
\Img\partial_{k+1}$. The three subspaces are mutually orthogonal because
$\partial_k\partial_{k+1} = 0$ (\Cref{thm:gauge:partial-squared}), and $\ker L_k$
is the orthogonal complement of $\Img(\partial_{k+1}) \oplus
\Img(\partial_k^{\!\top})$ in $C_k$.
\end{proof}

The dimension $\beta_k = \dim H_k(\mathcal{K}) = \operatorname{rank}(Z_k) -
\operatorname{rank}(B_k)$ is the $k$-th \emph{Betti number}, counting
$k$-dimensional holes: $\beta_0$ connected components, $\beta_1$ independent
loops, $\beta_2$ cavities. For the sphere $(\beta_0, \beta_1, \beta_2) = (1, 0,
1)$ and for the torus $(1, 2, 1)$, giving Euler characteristics $\chi(S^2) = 2$
and $\chi(T^2) = 0$ through the alternating sum $\chi = \sum_k (-1)^k \beta_k$.
On smooth manifolds the same invariants arise analytically: the \emph{de Rham
cohomology} $H^k_\dR(\Mani) = \ker(\mathrm{d}\colon \Omega^k \to \Omega^{k+1}) /
\Img(\mathrm{d}\colon \Omega^{k-1} \to \Omega^k)$ is isomorphic to singular
cohomology with real coefficients (de Rham's theorem), tying the exterior
calculus of \Cref{def:gauge:curvature} to topology. The exterior algebra
$\bigwedge V = \bigoplus_k \bigwedge^k V$, with antisymmetric product $v \wedge w
= -w \wedge v$, is precisely $\Cl(V, 0)$, the Clifford algebra of the zero form;
this is the bridge between the algebraic and topological extensions.

\subsection{Simplicial message passing and attention}
\label{sec:gauge:simplicial-attention}

The boundary structure dictates how cells exchange messages. For a $k$-simplex
$\sigma$ there are three natural neighbourhoods: the \emph{lower} neighbours (its
$(k-1)$-faces, via $\partial_k$), the \emph{upper} neighbours (the $(k+1)$-cells
of which it is a face, via $\delta_k$), and the \emph{horizontal} neighbours
(other $k$-cells sharing a face). \Cref{fig:gauge:simplicial-messages} shows the
three message types reaching a target edge.

\begin{figure}[ht]
\centering
\begin{tikzpicture}[>=Stealth, font=\small, scale=0.78, transform shape,
  msg/.style={line width=1.6pt, -{Stealth[length=6pt, width=4pt]}}]
\coordinate (A) at (0,0); \coordinate (B) at (5,0); \coordinate (C) at (2.5,3.5);
\coordinate (D) at (-3.5,1.2); \coordinate (E) at (-3.0,-2.0); \coordinate (F) at (8.5,1.2);
\fill[gdlorange!15] (A)--(B)--(C)--cycle;
\draw[gdlorange!40] (A)--(B)--(C)--cycle;
\node[gdlorange!70!black] at (2.5,1.6) {$\tau^2$};
\draw[gdlgray!50, line width=1.2pt] (A)--(C) (B)--(C);
\draw[gdlpurple!30, line width=2.2pt] (A)--(D) (A)--(E) (B)--(F);
\draw[gdlred!70, line width=4pt] (A)--(B);
\node[gdlred!80!black, font=\bfseries] at (2.5,-0.55) {$\sigma^1$ (target)};
\foreach \v/\lab/\pos in {A/$v_0$/{below left}, B/$v_1$/{below right}, C/$v_2$/above,
  D/$v_3$/left, E/$v_4$/{below left}, F/$v_5$/right} {
  \fill[gdlblue!60] (\v) circle (5pt);
  \node[font=\footnotesize, gdlblue!70!black, \pos=4pt] at (\v) {\lab};}
\draw[msg, gdlgreen!70] ($(A)+(-0.3,-1.8)$) to[bend right=25] ($(A)+(0.5,-0.25)$);
\draw[msg, gdlgreen!70] ($(B)+(0.3,-1.8)$) to[bend left=25] ($(B)+(-0.5,-0.25)$);
\node[gdlgreen!70!black, font=\scriptsize] at (2.5,-1.7) {$m^{\downarrow}$ (boundary)};
\draw[msg, gdlorange!70] (2.5,2.8) -- (2.5,0.3);
\node[gdlorange!70!black, font=\scriptsize, anchor=south west] at (3.0,2.3) {$m^{\uparrow}$ (coboundary)};
\draw[msg, gdlpurple!70] ($(A)!0.5!(D)$) to[bend right=20] ($(A)+(0.6,0.35)$);
\draw[msg, gdlpurple!70] ($(A)!0.5!(E)$) to[bend left=20] ($(A)+(0.6,-0.35)$);
\draw[msg, gdlpurple!70] ($(B)!0.5!(F)$) to[bend left=20] ($(B)+(-0.6,0.35)$);
\node[gdlpurple!70!black, font=\scriptsize, anchor=east] at (-3.8,0) {$m^{\leftrightarrow}$ (adjacent)};
\end{tikzpicture}
\caption[Simplicial messages]{Three message types reaching a target edge
$\sigma^1$ (red): \emph{lower} messages from its boundary vertices (green, via
$\partial$), \emph{upper} messages from triangles of which it is a face (orange,
via $\delta$), and \emph{horizontal} messages from adjacent edges (purple).}
\label{fig:gauge:simplicial-messages}
\end{figure}

\begin{definition}[Simplicial attention]
\label{def:gauge:simplicial-attention}
A simplicial attention network~\citep{giusti2022simplicial} updates the
representation $\mathbf{h}_\sigma$ of a $k$-simplex by aggregating the three
message types,
\[
  \mathbf{h}_\sigma^{(\ell+1)} = \sigma\!\Big(\mathbf{W}^{\text{self}}
  \mathbf{h}_\sigma^{(\ell)} + \mathbf{m}_\sigma^{\downarrow} +
  \mathbf{m}_\sigma^{\uparrow} + \mathbf{m}_\sigma^{\leftrightarrow}\Big),
  \qquad
  \mathbf{m}_\sigma^{\bullet} = \sum_\tau \alpha_{\tau\to\sigma}\,
  \mathbf{W}^{\bullet}\mathbf{h}_\tau,
\]
with attention coefficients
\[
  \alpha_{\tau\to\sigma} = \softmax_\tau\Big(\mathrm{LeakyReLU}\big(\mathbf{a}^{\!\top}
  [\mathbf{W}_\sigma\mathbf{h}_\sigma \,\|\, \mathbf{W}_\tau\mathbf{h}_\tau \,\|\,
  \mathbf{e}_{\text{orient}}(\tau, \sigma)]\big)\Big)
\]
that incorporate the relative orientation $\mathbf{e}_{\text{orient}}$. Folding a
sign $s_{\tau\to\sigma} \in \{\pm 1\}$ into the coefficients yields
\emph{oriented} (signed) attention, preserving cohomological structure.
\end{definition}

\begin{theorem}[Simplicial equivariance]
\label{thm:gauge:simplicial-equivariance}
Simplicial attention networks are equivariant under simplicial automorphisms: if
$\phi\colon \mathcal{K} \to \mathcal{K}$ is an automorphism, then
$\mathrm{SAN}(\phi(\mathbf{X})) = \phi(\mathrm{SAN}(\mathbf{X}))$.
\end{theorem}

\begin{proof}
An automorphism $\phi$ induces a permutation $P_k$ of the $k$-simplices
preserving incidence, $B_k P_{k-1} = P_k B_k$, hence commuting with the
Laplacians, $P_k L_k P_k^{\!\top} = L_k$. The update depends on the cells only
through these incidence matrices and through attention coefficients that are
symmetric functions of the representations, so $P_k\cdot\mathrm{SAN}(\mathbf{X}) =
\mathrm{SAN}(P_k\cdot\mathbf{X})$ in each dimension $k$.
\end{proof}

The Hodge Laplacian also serves as a spectral filter for attention. With the
eigendecomposition $L_k = U_k\Lambda_k U_k^{\!\top}$ (multiplicity of the zero
eigenvalue equal to $\beta_k$, with harmonic eigenvectors), one builds queries,
keys, and values from spectral filters $\mathbf{g}_\theta(L_k) = \sum_i \theta_i
T_i(\tilde L_k)$ in the Chebyshev basis, giving \emph{Hodge Laplacian
attention} that separates irrotational, solenoidal, and harmonic components of
edge signals. \emph{Persistent homology} adds a multiscale dimension: over a
filtration $\mathcal{K}_0 \subseteq \cdots \subseteq \mathcal{K}_m$ the
persistence diagram records the birth $b$ and death $d$ of each homological
feature, and persistence values $\mathrm{pers}(i, j)$ can bias attention
logits, $\alpha_{ij} = \softmax(\mathbf{q}_i^{\!\top}\mathbf{k}_j / \sqrt{d} + w\,
\mathrm{pers}(i,j))$, lending robustness to noise and a connection to the
topology of the data. Finally, \emph{cellular sheaves} attach a vector space
(stalk) to each cell and a restriction map to each incidence, and the sheaf
Laplacian $L_\mathcal{F} = \delta^{\!\top}\delta$ generalizes the graph Laplacian
to heterogeneous feature spaces, underpinning sheaf neural
networks~\citep{hansen2021sheaf, bodnar2023cellular}. Across all these
constructions the message-passing principle recurs, now respecting incidence and
orientation rather than mere adjacency.

\section{Positional encoding, symmetry breaking, and quantum extensions}
\label{sec:gauge:positional}

Two further ingredients refine geometric architectures: how to encode position
on a domain, and how to relax equivariance when the data demand it.

\paragraph{Fourier features and Bochner's theorem.}
A translation-invariant positive-definite kernel admits a spectral
representation, the foundation of random Fourier features.

\begin{theorem}[Bochner]
\label{thm:gauge:bochner}
A continuous kernel $k(\mathbf{x}, \mathbf{y}) = k(\mathbf{x} - \mathbf{y})$ on
$\mathbb{R}^d$ is positive definite if and only if it is the Fourier transform of
a finite non-negative measure $\mu$:
\[
  k(\mathbf{x} - \mathbf{y}) = \int_{\mathbb{R}^d} e^{2\pi i\,
  \omega^{\!\top}(\mathbf{x} - \mathbf{y})}\, \mathrm{d}\mu(\omega).
\]
\end{theorem}

Sampling frequencies $\omega_j$ from the associated density yields the feature
map $z(\mathbf{x}) = \sqrt{2/m}\,[\cos(2\pi\omega_j^{\!\top}\mathbf{x}),
\sin(2\pi\omega_j^{\!\top}\mathbf{x})]_j$ with $\mathbb{E}[z(\mathbf{x})^{\!\top}
z(\mathbf{y})] = k(\mathbf{x}, \mathbf{y})$ and uniform error $O(1/\sqrt{m})$. On
a compact Riemannian manifold the sinusoids are replaced by Laplace--Beltrami
eigenfunctions $\{\phi_k\}$: the map $\gamma_\Mani(p) = \sqrt{2/m}\,(\phi_1(p),
\ldots, \phi_m(p))$ approximates the diffusion kernel $k_t(p, q) = \sum_k
e^{-t\lambda_k}\phi_k(p)\phi_k(q)$ with error $O(e^{-t\lambda_{m+1}})$, the
intrinsic analogue of Fourier features. In a principal bundle these can be made
gauge-equivariant by transporting them, $\gamma_{\text{gauge}}(p) =
\Gamma_\gamma^{-1}\gamma(\pi(p))$, so that the encoding transforms correctly under
local frame changes.

\paragraph{Controlled symmetry breaking.}
Exact equivariance is sometimes too rigid: real symmetries may be approximate, or
privileged directions may carry information. \citet{vanderouderaa2024symmetry}
interpolate between an equivariant map $\mathcal{F}_{\text{eq}}$ and an
unconstrained one,
\[
  \mathcal{F}_\lambda = (1 - \lambda)\,\mathcal{F}_{\text{eq}} + \lambda\,
  \mathcal{F}_{\text{in}},
  \qquad \lambda \in [0, 1].
\]
At $\lambda = 0$ equivariance is exact; as $\lambda$ grows a controlled
directional bias is introduced, with $\lambda$ learned and regularized toward
zero so the network defaults to equivariance unless the data justify breaking it.

\begin{proposition}[Smooth transition]
\label{prop:gauge:symmetry-breaking}
If $f_0$ is $G$-equivariant and $h$ arbitrary, then $f_\lambda = (1 - \lambda)
f_0 + \lambda h$ satisfies $\|f_\lambda(g\cdot x) - \rho(g) f_\lambda(x)\| =
\lambda\,\|h(g\cdot x) - \rho(g) h(x)\|$, which tends to $0$ uniformly on compacts
as $\lambda \to 0$.
\end{proposition}

\begin{proof}
Expanding and using the exact equivariance of $f_0$, the $f_0$ terms cancel,
leaving $\lambda\,(h(g\cdot x) - \rho(g) h(x))$. On a compact $K \subset \Mani
\times G$ the continuous function $(x, g) \mapsto \|h(g\cdot x) - \rho(g) h(x)\|$
is bounded by some $M$, so the deviation is at most $\lambda M \to 0$.
\end{proof}

These principles extend even to quantum circuits, where a gauge-equivariant
quantum network is a parametrized unitary $U(\theta) = \exp(i\sum_j \theta_j
H_j^{\text{gauge}})$ built from Hamiltonians commuting with the gauge generators,
mirroring the classical construction at the level of unitary evolution.

\section{Synthesis: gauges in the unified picture}
\label{sec:gauge:synthesis}

The gauge framework completes the five-G arc, and the same message-passing thread
that ran through grids, groups, graphs, and geodesics runs through it. A
geometric convolution is a structured message-passing layer: the underlying graph
reflects the geometry of the domain, the message functions preserve equivariance
to the symmetry group, aggregation respects local geometry (convolution becomes a
locally weighted sum), and weights are shared across the orbits of the group.
Gauge-equivariant convolution is this principle on a principal bundle, with
parallel transport supplying the message functions that align neighbouring
frames.

Seen this way the architectures of the whole book are facets of one construction.
Classical and group convolutions are the flat, globally trivialized case;
graph and geodesic message passing relax the domain to a discrete or curved one;
gauge convolution adds the local frame and its connection; and attention ---
including the Clifford and simplicial variants of this chapter --- is message
passing with adaptive, content-dependent weights. The transition is even
quantitative: a geometric convolution arises as the sharp-temperature limit of an
attention kernel,
\[
  (f * g)(x) = \lim_{\beta \to \infty} \int_\Mani
  \frac{e^{-\beta d(x, y)^2}}{\int_\Mani e^{-\beta d(x, y')^2}\, \mathrm{d}y'}\,
  g(y)\, f(y)\, \mathrm{d}y,
\]
recovering a localized convolution from soft attention as $\beta \to \infty$.
\Cref{ch:applications} returns to the concrete payoffs --- meshes, spheres, and
$3$D molecules --- and \Cref{ch:conclusion} draws together the unified
message-passing view of the entire book, of which the gauge picture is the most
general expression.
